\section{The stationary critical exponent}
Write
\[
 E(a)=\frac{\int_0^{2\pi}\sin^2\theta\,e^{a\cos\theta}\,d\theta}
             {\int_0^{2\pi}e^{a\cos\theta}\,d\theta}.
\]
This is \texttt{vonMisesSRatio}. Its denominator is strictly positive.

\subsection{The second-order expansion}
\begin{claim}
As $a\to0$,
\[
 E(a)=\frac12-\frac{a^2}{16}+o(a^2).
\]
\end{claim}
\begin{proof}
The previously proved differentiability results give $E\in C^2$,
$E(0)=1/2$, $E'(0)=0$, and $E''(0)=-1/8$. Taylor's theorem with a little-o
remainder yields
$E(a)=E(0)+E'(0)a+E''(0)a^2/2+o(a^2)$.
Substituting those values gives the claimed polynomial. The Lean proof applies
the general Taylor theorem, identifies its second iterated derivative using
the derivative lemmas, and simplifies its Taylor polynomial.
\end{proof}
\formal{second-order}
\scope The asserted remainder is $o(a^2)$. This statement supplies no
fourth-order remainder estimate.

\subsection{An arbitrary family approaching threshold}
\begin{claim}
Let $D>0$, let $\ell$ be a filter on an index type, and let $K_x,r_x$ be
real-valued functions on that type. Suppose $K_x\to2D$, $r_x\to0$ along
$\ell$, and eventually
\[
 K_x>2D,\qquad r_x>0,\qquad
 r_x=\operatorname{selfConsistency}(K_x,D,r_x).
\]
Then $r_x^2/(K_x-2D)\to1/D$ along $\ell$.
\end{claim}
\begin{proof}
Set $a_x=K_xr_x/D$ and write
\[
 E(a)=\frac12-\frac{a^2}{16}+a^2q(a).
\]
The previous expansion gives $q(a)\to0$ as $a\to0$, taking the quotient
with Lean's totalised division at zero. The assumed limits imply $a_x\to0$,
hence $q(a_x)\to0$.

For an index satisfying the eventual hypotheses, positivity gives
$a_x\ne0$ and $K_x-2D\ne0$. The previously proved reformulation of the
self-consistency equation gives $E(a_x)=D/K_x$. Substituting the expansion
and rearranging yields
\[
 K_x-2D=K_xa_x^2\bigl(1/8-2q(a_x)\bigr),\qquad
 \frac{r_x^2}{K_x-2D}
 =\frac{D^2}{K_x^3\bigl(1/8-2q(a_x)\bigr)}.
\]
The denominator in the last expression tends to
$(2D)^3/8=D^3>0$. The quotient therefore tends to $D^2/D^3=1/D$.
The equality of the quotients holds eventually, which transfers this limit
to the desired expression.
\end{proof}
\formal{family-exponent}
\scope Convergence to zero and eventual stationarity are hypotheses. This
general statement allows the bottom filter, where its eventual and convergence
conditions are vacuous. A concrete branch supplies the substantive application
below. The claim establishes stationary asymptotics, with no time evolution
or stability conclusion.

\subsection{A branch that satisfies the hypotheses}
For $D>0$ and $K>2D$, define $r_*(D,K)$ by choosing the coherent solution from
the theorem that proves existence and uniqueness in $(0,1]$.
The Lean function \texttt{coherentBranch} uses this choice in that regime
and returns zero outside it.

\begin{claim}
If $D>0$ and $K>2D$, then $0<r_*(D,K)\leq1$ and
\[
r_*(D,K)=\operatorname{selfConsistency}(K,D,r_*(D,K)).
\]
\end{claim}
\begin{proof}
The two inequalities on $D,K$ select the positive branch of the definition.
The specification of the chosen witness is exactly positivity, the bound by
one, and the fixed-point equation supplied by the existence theorem.
\end{proof}
\formal{branch-spec}

\begin{claim}
For fixed $D>0$, as $K$ decreases to $2D$ through $K>2D$,
\[
 \frac{r_*(D,K)^2}{K-2D}\longrightarrow\frac1D.
\]
\end{claim}
\begin{proof}
Apply the family theorem to the right-hand neighbourhood filter of $2D$,
with the identity function for $K$. The previously proved
\texttt{coherentBranch\_tendsto\_zero} supplies the limit of $r_*$ at the
threshold. Membership in the right-hand neighbourhood gives $K>2D$
eventually, and the branch specification supplies positivity and stationarity
there. These are all the hypotheses of the family theorem.
\end{proof}
\formal{branch-exponent}
\scope The approach filter on the real line is nontrivial, and the branch
provides actual coherent solutions. This removes the vacuity issue in this
application. The conclusion still concerns stationary solutions.
